\documentclass[12pt, aspectratio=169]{beamer}
\usepackage[utf8]{inputenc}
\usepackage[portuguese]{babel}
\usepackage{graphicx}
\usepackage{xcolor}

% Marketing Intelligence Intelligence - Paleta Oficial
\definecolor{Marketing IntelligenceObsidian}{HTML}{0A0C0F}
\definecolor{Marketing IntelligencePergaminho}{HTML}{F4F2ED}
\definecolor{Marketing IntelligenceVerdeSinal}{HTML}{C8F542}
\definecolor{Marketing IntelligenceNevoa}{HTML}{8A8F99}

\setbeamercolor{background canvas}{bg=Marketing IntelligenceObsidian}
\setbeamercolor{title}{fg=Marketing IntelligenceVerdeSinal}
\setbeamercolor{author}{fg=Marketing IntelligenceNevoa}
\setbeamercolor{date}{fg=Marketing IntelligenceNevoa}
\setbeamercolor{frametitle}{fg=Marketing IntelligenceVerdeSinal}
\setbeamercolor{normal text}{fg=Marketing IntelligencePergaminho}
\setbeamercolor{block title}{bg=Marketing IntelligenceVerdeSinal, fg=Marketing IntelligenceObsidian}
\setbeamercolor{itemize item}{fg=Marketing IntelligenceVerdeSinal}

\title{Marketing Intelligence MI --- Startup Signal Intelligence}
\subtitle{Marketing Intelligence \& Reputation Analysis}
\author{Renato Assis $|$ Marketing Intelligence Intelligence}
\institute{Brasil --- UFPB (Data Science for Business)}
\date{Março, 2026}

\begin{document}

\begin{frame}
  \titlepage
\end{frame}

\begin{frame}{Contexto Estratégico: O Problema}
  \begin{itemize}
    \item \textbf{Maturidade Analítica (Davenport):} Muitas organizações estão no \textit{Estágio 1} (Analiticamente Deficientes).
    \item \textbf{Decisões Baseadas em Intuição:} A carência de inteligência em dados não estruturados gera crises de reputação.
    \item \textbf{Oportunidade:} Transformar narrativas digitais em sinais preditivos de sucesso e tração de mercado.
  \end{itemize}
\end{frame}

\begin{frame}{Marketing Intelligence MI: A Solução}
  \begin{block}{Objetivo do Módulo MI}
    Implementar um sistema de \textbf{Signal Intelligence} capaz de analisar a saúde comunicacional de startups e marcas a partir de textos públicos (GitHub, Twitter, Blogs).
  \end{block}
  \vspace{0.3cm}
  \centering
  \textbf{Fluxo:} Coleta $\rightarrow$ Limpeza (RE) $\rightarrow$ NLP (spaCy/NLTK) $\rightarrow$ Inteligência (Transformers)
\end{frame}

\begin{frame}{Etapa 1: Estruturação Narrativa (RE + spaCy)}
  \textbf{Técnica: Expressões Regulares (Regex)}
  \begin{itemize}
    \item \textbf{Objetivo:} Extrair "Trust Markers" e limpar ruído.
    \item \textbf{O que buscamos:} E-mails, versões de software, menções, referências de issues e métricas de qualidade de README.
    \item \textbf{Significado:} É a base do \textit{Narrative Identity Card}, transformando texto solto em dados estruturados.
  \end{itemize}
  \vspace{0.2cm}
  \textbf{Técnica: spaCy Identity Setup}
  \begin{itemize}
    \item Reconhecimento de Entidades Nomeadas (NER) e Extração de Propostas de Valor.
  \end{itemize}
\end{frame}

\begin{frame}{Fundamentação Científica}
  \small
  \begin{itemize}
    \item \textbf{Qiu et al. (2025):} BERTweet para predição de financiamento.
    \item \textbf{Peixoto et al. (2023):} LDA para mapear fases de vida de startups.
    \item \textbf{Saura et al. (2019):} Indicadores de sucesso via Machine Learning.
    \item \textbf{Jin, Wu \& Hitt (2017):} Sinais sociais como proxies de captação futura.
    \item \textbf{Cheng (2024):} O impacto da presença multicanal na legitimidade ($R^2 \uparrow$).
  \end{itemize}
\end{frame}

\begin{frame}{Status Atual do Projeto}
  \centering
  \begin{itemize}
    \item[\checkmark] Estrutura de Diretórios e README Finalizada.
    \item[\checkmark] Identidade Visual Definida (Signal Green/Obsidian).
    \item[\checkmark] Módulo \texttt{src/limpeza.py} Concluído (Regex Integrado).
    \item[$\rightarrow$] Próximo: Motor de Coleta de Dados (API/Scraping).
  \end{itemize}
\end{frame}

\end{document}
